\section{The first wild index for every prime at least five}\label{sec:first-wild-general}
This section proves the exact index at $e=p(p-1)$ by combining the graph estimate with an actual cyclic-point construction.

\begin{proposition}\label{prop:first-wild-ss}
Let $p\ge5$, $\mu_p\subset K$ and $e_K=p(p-1)$, with $K$ complete and its residue field perfect. Two whole potentially supersingular height-two groups with isomorphic $p^3$-torsion have equal special inertia characters, regardless of their residual inertia type.
\end{proposition}
\begin{proof}
If the characters differed, Proposition~\ref{fc-thm:valuation}, which has no condition on $v_p(e_K)$, would imply
\[
                    p(p-1)(p^3-1)\ge(p-1)p^4.
\]
The left side minus the right side equals $-p(p-1)<0$. Equivalently, the primitive-coordinate average $(p-1)/p^3$ is strictly below the required minimum $(p-1)/(p^3-1)$. This excludes every residual branch.
\end{proof}

\begin{proof}[Uniform upper bound in Theorem~\ref{thm:first-wild-general}]
Pass to the completed maximal unramified base, preserving $e$ and inertia. For an allowed height-two group the determinant on inertia is $\chi_p^d\epsilon$, where $0\le d\le2$ is its total dimension and $\epsilon$ is a finite $\Z_p$-valued character of order dividing $p-1$, as in Theorem~\ref{thm:second-wild-exact}.

The cyclotomic image modulo $p^3$ has nontrivial wild part. To see this without assuming $\mu_{p^2}\not\subset K$, take a Cohen coefficient field $K_0$. The polynomial $\Phi_{p^3}(1+X)$ is monic, has constant term $p$ and reduction $X^{p^2(p-1)}$, so is Eisenstein over $\Ocal_{K_0}$. The extension $K_0(\mu_{p^3})/K_0$ is totally ramified of index $p^2(p-1)$, which cannot divide $e_K=p(p-1)$. Thus $\mu_{p^3}\not\subset K$. Since $\mu_p\subset K$, the remaining cyclotomic image is a nontrivial $p$-group.

An isomorphism of $p^3$-torsion identifies determinants modulo $p^3$. On wild inertia the finite factors are trivial, so $p\mid d-d'$. As $|d-d'|\le2<p$, one has $d=d'$. Reduction modulo $p$ then identifies $\epsilon,\epsilon'$, because reduction is injective on roots of unity of order dividing $p-1$.

For $d=0$, all potentially good graded pieces are potentially \'etale; their special character is the Brauer lift of the residual semisimplification. Hence it is determined by the common $p$-torsion. Cartier duality gives $d=2$. If $d=1$, the group either has two height-one graded pieces or is a whole potentially good height-two group. In the latter case ordinary reduction gives the connected--\'etale filtration, which descends by uniqueness. The remaining whole possibility is supersingular.

Two height-one-filtered groups are compared residually, since $\chi_p\equiv1\pmod p$ over $K$. A mixed comparison is impossible: Proposition~\ref{prop:mixed-obstruction}, with $n=3$ and $v_p(e)=1$, would require $e>p^2$, contrary to $e=p^2-p$. Two whole supersingular groups are covered by Proposition~\ref{prop:first-wild-ss}. All dimensions and filtration types are accounted for, proving the uniform upper bound two.
\end{proof}

\begin{proof}[Actual lower bound in Theorem~\ref{thm:first-wild-general}]
We specialize Theorem~\ref{thm:general-arc}, whose hypotheses do not require $p\nmid e$, and give the parameters and exact point construction. Put
\[
 k=\overline{\F}_p,\quad W=W(k),\quad q=p^2,\quad N=q-1,
 \qquad d=Np(p-1),\qquad b=p(p-1)(p^2-p-1).
\]
Use the integral universal group $F_u$ and its homogeneous logarithm from \eqref{eq:homogeneous-logarithm}. Define
\[
 Y_0=X,\qquad Y_1=[p]_{X^b}(X),\qquad
 P(X)=\Phi_{X^b}(Y_1),\qquad \Phi_u(T)=[p]_u(T)/T.
\]
The two strict leading-term comparisons are
\[
 b-(q-p)=Np(p-2)>0,\qquad Nq-d=Np>0.
\]
Modulo $p$, the first selects $T^q$ in $Y_1$, so $Y_1=X^q+O(X^{q+1})$. In $P$ the term $uY_1^{p-1}$ has degree $b+(p-1)q=d$, strictly before $Y_1^{q-1}$. Intervening coefficients are divisible by $u$ modulo $p$ and are later; higher coefficients are integral. Consequently
\[
                     P(0)=p,\qquad P(X)=X^d+O(X^{d+1})\pmod p.
\]
Weierstrass preparation gives an Eisenstein polynomial of degree $d$, so $R=W[[X]]/(P)$ is a complete discrete valuation ring with uniformizer $x$, residue field $k$ and $v_R(p)=d$. Since $q<d$, the value of $Y_1(x)$ is exactly $q$, hence is nonzero. The relation $\Phi(Y_1(x))=0$ makes $[p]Y_1(x)=0$. Thus $x$ has exact order $p^2$ in $F_{x^b}$.

Choose $\zeta\in W$ of order $N$ and let $a=\zeta^{-1}$. The congruence $b\equiv-(p-1)\pmod N$ and the exact homogeneous identities give
\[
 [p]_{a^{p-1}u}(a^{-1}T)=a^{-1}[p]_u(T),\qquad
 \Phi_{a^{p-1}u}(a^{-1}T)=\Phi_u(T).
\]
Therefore $X\mapsto\zeta X$ fixes $P$ and defines an automorphism $\sigma$ of $R$ of exact order $N$. With $L=\operatorname{Frac}R$ and $K=L^{\langle\sigma\rangle}$, the extension $L/K$ is totally tame of degree $N$, so $e_K=d/N=p(p-1)$ and the residue field is $k$. The full formal-group isomorphism $\sigma^*F\to F$, $T\mapsto aT$, follows from the homogeneous logarithm, satisfies the cocycle condition and fixes the point because $a\sigma(x)=x$. Effective descent of the generic finite levels gives an actual height-two group $G/K$ with a rational point of exact order $p^2$ and special eigenvalues $a,a^p$.

The field contains $\mu_p$. Indeed $x^N$ is a uniformizer of $K$; write $-p=(x^N)^eU$ with $U$ a unit. Hensel's lemma provides a $(p-1)$st root of $U$, and $(p-1)\mid e$. Thus $K$ contains a $(p-1)$st root of $-p$, generating the unique tame degree-$(p-1)$ extension of its Cohen field, namely the cyclotomic extension.

Finally apply the actual point-to-ambiguity construction of Lemmas~\ref{lem:point-ambiguity} and~\ref{lem:extension-lift}. More explicitly, write $\det G=\chi_p\psi$, let $C$ be the constant cyclic subgroup of order $p^2$ generated by the point, and form $G'=G/C$. The isogeny sequence gives
\[
0\longrightarrow(\Z/p^2)(\chi_p\psi)
 \longrightarrow G'[p^2]\longrightarrow\Z/p^2\longrightarrow0.
\]
Twisting by $\psi^{-1}$ gives a whole potentially supersingular group $M$ and an extension of $(\Z/p^2)(\psi^{-1})$ by $\mu_{p^2}$. The extension-lifting lemma realizes this finite extension as $H[p^2]$ for an actual height-two group $H$ with graded pieces $\mu_{p^\infty}$ and $(\Q_p/\Z_p)(\psi^{-1})$.

On a tame generator, the special eigenvalues of $M$ are $a^{-1},a^{-p}$. Those of $H$ are $1,a^{-(p+1)}$. Their trace difference is
\[
                -(a^{-1}-1)(a^{-p}-1)\ne0.
\]
Hence $M[p^2]\simeq H[p^2]$ with unequal special characters at the required absolute index. This proves the lower bound two and completes Theorem~\ref{thm:first-wild-general}.
\end{proof}
